% !TEX TS-program = pdflatex
% !TEX encoding = UTF-8 Unicode
% Numerical methods: a page of the wiki "Enhanced dissipation for passive scalars", generated from the wiki on 10 November 2026.
% It compiles as it is with pdflatex (Overleaf, TeXShop, TeXstudio, ...); typeset twice for the citations.
% To propose a change: keep the preamble and the line \title{...}; keep only the blocks you change
% (a \section{...} with its content, or a statement or proof with its \label{...}), change them, and post
% the file in a new topic of the Zulip channel of the part. The rules: page "How to contribute" of the wiki.
\documentclass[11pt]{article}
\usepackage[utf8]{inputenc}
\usepackage[T1]{fontenc}
\usepackage{lmodern}
\usepackage[a4paper,margin=2.5cm]{geometry}
\usepackage{amsmath,amssymb,amsthm}
\usepackage{graphicx,tabularx,listings}
\usepackage[skins,breakable]{tcolorbox}
\usepackage[round]{natbib}
\usepackage[hidelinks]{hyperref}
\setlength{\emergencystretch}{2em}

\newtheorem{theorem}{Theorem}
\newtheorem{proposition}[theorem]{Proposition}
\newtheorem{lemma}[theorem]{Lemma}
\newtheorem{corollary}[theorem]{Corollary}
\newtheorem{conjecture}[theorem]{Conjecture}
\newtheorem*{claim}{Claim}
\theoremstyle{definition}
\newtheorem{definition}[theorem]{Definition}
\newtheorem{example}[theorem]{Example}
\theoremstyle{remark}
\newtheorem{remark}[theorem]{Remark}

%% The commands of the wiki. This preamble is the same for every page: do not change it, and do not
%% define commands of your own (the wiki has none).
\makeatletter
% after a statement: \begin{idea}, \begin{architecture}, \begin{note}[Title] (and \begin{proof}[Title])
\newcommand{\ed@runin}[1]{\par\topsep6\p@\@plus6\p@\relax\trivlist\item[\hskip\labelsep\itshape#1\@addpunct{.}]\ignorespaces}
\newenvironment{idea}[1][Idea]{\ed@runin{#1}}{\endtrivlist\@endpefalse}
\newenvironment{architecture}[1][Proof architecture]{\ed@runin{#1}}{\endtrivlist\@endpefalse}
% \begin{note}[Title]: an aside, in a dashed box; its title in italics at the start of the text
\tcbset{ed@note/.style={enhanced jigsaw,breakable,frame hidden,interior hidden,arc=2pt,borderline={0.5pt}{0pt}{black!50,dash pattern=on 3pt off 2pt},
  left=3mm,right=3mm,top=2mm,bottom=2mm,before skip=3mm plus 1mm,after skip=3mm plus 1mm}}
\newenvironment{note}[1][Note]{\begin{tcolorbox}[ed@note]\noindent{\itshape#1\@addpunct{.}}\enspace\ignorespaces}{\end{tcolorbox}}
% text still to be written; labels of the text; the Lean banner of the statement above
\newcommand{\placeholder}[1]{\if@inlabel\else\par\noindent\fi\textit{[#1]}\@ifnextchar\end{}{\par}}
\newcommand{\ed@label}[1]{\mbox{\footnotesize[#1]}}
\newcommand{\demo}[1]{\ed@label{demo text}}
\newcommand{\heuristic}[1]{\ed@label{heuristic}}
\newcommand{\numerics}[1]{\ed@label{numerics}}
\newcommand{\published}[1]{\ed@label{published}}
\newcommand{\preprint}[1]{\ed@label{preprint, to be verified}}
\newcommand{\raw}[1]{\ed@label{raw, unchecked}}
\def\ed@formalized{formalized}\def\ed@image{image}\def\ed@video{video}\def\ed@diagram{diagram}
\newcommand{\lean}[2][]{\par{\noindent\footnotesize\def\ed@a{#2}Lean: \ifx\ed@a\ed@formalized formalized\ifx\relax#1\relax\else\
  (\href{#1}{certificate})\fi\else not yet formalized\fi.\par}}
% the contributors of a section, most recent first (chosen by the editors), in small type under its title
\newcommand{\contributors}[1]{\par{\noindent\footnotesize\itshape Contributors: #1\par}\smallskip}
% the two lines under the title: the page's team and status, its Zulip topic (default: the title)
\def\ed@title{}\let\ed@maketitle\maketitle\renewcommand{\maketitle}{\let\ed@title\@title\ed@maketitle}
\newcommand{\pagestatus}[1]{{\par\centering\small#1\par}}
\newcommand{\zulip}[2][]{{\par\centering\small Discuss on Zulip: \##2, topic
  ``\ifx\relax#1\relax\ed@title\else#1\fi''.\par}\medskip}
% links to other pages of the wiki, anchors, wiki code shown as text
\newcommand{\wikilink}[2]{\ifx\relax#2\relax#1\else#2\fi}
\newcommand{\wikianchor}[1]{\phantomsection\label{#1}}
\newcommand{\wikipage}[1]{}% the first line of a page file of the whole-wiki project (pages/<page>.tex): the page's title
\newcommand{\nowikicode}[1]{\texttt{\frenchspacing#1}}
\newcommand{\nowiki}[1]{#1}
\newcommand{\wikitemplate}[1]{\texttt{\{\{#1\}\}}}
\newcommand{\wikidiv}[2]{\par#2\par}
% \begin{quote}: a sentence set apart, such as the central question of the Overview; centred, in italics
\renewenvironment{quote}{\list{}{\rightmargin\leftmargin}\item\relax\centering\itshape}{\endlist}
% pictures, videos, diagrams: a grey panel with the picture and its caption. The picture (for a video, its poster image)
% is shown when its file, or a .pdf or .png file of the same name, is next to this file, in figures/ or in ../media
\tcbset{ed@media/.style={enhanced,colback=black!4,frame hidden,arc=2pt,left=3mm,right=3mm,top=2.5mm,bottom=2mm,
  before skip=3mm plus 1mm,after skip=3mm plus 1mm}}
\def\ed@poster{}\def\ed@record{}\def\ed@svg{svg}
\define@key{ed@media}{poster}{\def\ed@poster{#1}}\define@key{ed@media}{record}{\def\ed@record{#1}}\define@key{ed@media}{id}{}
\newcommand{\ed@find}[1]{\IfFileExists{#1}{\edef\ed@found{#1}}{\IfFileExists{figures/#1}{\edef\ed@found{figures/#1}}{%
  \IfFileExists{../media/#1}{\edef\ed@found{../media/#1}}{}}}}
\newcommand{\ed@picturefile}[1]{\def\ed@found{}\edef\ed@b{#1}\ifx\ed@b\@empty\else\filename@parse{\ed@b}\edef\ed@base{\filename@base}%
  \edef\ed@ext{\filename@ext}\ifx\ed@ext\ed@svg\else\ed@find{\ed@b}\fi
  \ifx\ed@found\@empty\ed@find{\ed@base.pdf}\fi\ifx\ed@found\@empty\ed@find{\ed@base.png}\fi\fi}
\newcommand{\ed@panel}[2]{\begin{tcolorbox}[ed@media,#1]\ifx\ed@found\@empty\else{\centering\includegraphics[width=\linewidth,
  height=0.35\textheight,keepaspectratio]{\ed@found}\par}\smallskip\fi{\small\textbf{\ifx\ed@a\ed@video Video\else\ifx\ed@a\ed@diagram
  Diagram\else Figure\fi\fi.}\enspace#2\par}{\footnotesize\itshape\ifx\ed@a\ed@video The video plays in the online edition of this page.\fi
  \ifx\ed@record\@empty\else\ifx\ed@a\ed@video\ \fi How it was made: page Numerical methods, record \texttt{\ed@record}.\fi\par}
  \end{tcolorbox}}
% in the text, the panel is a figure (it floats to where it fits); inside a box (a note, a table), it stays in place
\renewcommand{\topfraction}{0.85}\renewcommand{\bottomfraction}{0.6}\renewcommand{\textfraction}{0.1}\renewcommand{\floatpagefraction}{0.8}
\newcommand{\media}[4][]{\par\def\ed@poster{}\def\ed@record{}\setkeys{ed@media}{#1}\def\ed@a{#2}%
  \ifx\ed@a\ed@video\ed@picturefile{\ed@poster}\else\ed@picturefile{#3}\fi
  \ifinner\ed@panel{}{#4}\else\begin{figure}[htbp]\ed@panel{before skip=0pt,after skip=0pt}{#4}\end{figure}\fi}
% tables (\lv{label}{text}; \header{a & b}, \row{a & b}) and code
\newenvironment{levels}[1][]{\par\smallskip\noindent\tabularx{\linewidth}{@{}>{\bfseries\raggedright\arraybackslash}p{0.2\linewidth}X@{}}}{\endtabularx\par\smallskip}
\newcommand{\lv}[2]{#1 & #2\\[3pt]}
\newenvironment{wikitable}[1]{\par\smallskip\noindent\tabularx{\linewidth}{|*{#1}{>{\raggedright\arraybackslash}X|}}\hline}{\endtabularx\par\smallskip}
\newcommand{\header}[1]{#1\\\hline\hline}
\newcommand{\row}[1]{#1\\\hline}
\lstnewenvironment{wikicode}[1][]{\lstset{basicstyle=\ttfamily\footnotesize,breaklines=true,columns=fullflexible,
  keepspaces=true,showstringspaces=false,frame=single}\ifx\relax#1\relax\else\lstset{language=#1}\fi}{}
% the heading of the reference list is the section References of the page
\renewcommand{\bibsection}{}
% characters of the wiki's text that LaTeX does not know
\DeclareUnicodeCharacter{2192}{\ensuremath{\rightarrow}}\DeclareUnicodeCharacter{203A}{\ensuremath{\rangle}}
\DeclareUnicodeCharacter{2713}{\ensuremath{\checkmark}}\DeclareUnicodeCharacter{25CB}{\ensuremath{\circ}}
\DeclareUnicodeCharacter{2248}{\ensuremath{\approx}}\DeclareUnicodeCharacter{2264}{\ensuremath{\leq}}
\DeclareUnicodeCharacter{2265}{\ensuremath{\geq}}\DeclareUnicodeCharacter{2212}{\ensuremath{-}}
\makeatother

\begin{document}
\title{Numerical methods}
\author{Wiki \emph{Enhanced dissipation for passive scalars}}
\date{10 November 2026}
\maketitle
\pagestatus{Page team: to be named · Status: skeleton, one record}
\zulip{numerics}

Every computed picture and video of this wiki has a record on this page: what was computed, how, with which parameters and checks, and the code that produced it.

\section{How a record works}
\begin{itemize}
  \item A record is a section with an anchor \texttt{num-…}; the picture or video links to it with the parameter \texttt{record} of \wikilink{Template:Media}{Media}.
  \item It gives the equation, the parameters, the method, the checks, the code and how to run it, and its status.
  \item The code is in the folder \texttt{numerics/} of the \href{https://github.com/endipa/endipa.github.io}{repository} of the wiki, one folder per record, with the output of its checks.
\end{itemize}

\section{Records}
\subsection{Heat versus shear on the torus}\label{num-heat-vs-shear}
\begin{levels}
\lv{Produces}{the video on the \wikilink{Overview}{} (\texttt{heat-vs-shear.mp4}) and the check figure below}
\lv{Equation}{$\partial_t f+\sin y\,\partial_x f=\nu\Delta f$ on $[0,2\pi)^2$, periodic, and the heat equation $\partial_t f=\nu\Delta f$; same datum $f_0=\cos x$}
\lv{Parameters}{$\nu=10^{-3}$, times $0\le t\le 100$; video of 20 s at 24 frames per second, time running uniformly (5 time units per second), last frame held 1.5 s}
\lv{Method}{one Fourier mode in $x$: $f=\operatorname{Re}(g(t,y)e^{ix})$ with $\partial_t g=-i\sin y\,g+\nu(\partial_y^2-1)g$; 512 Fourier modes in $y$; Strang splitting with time step 0.01, transport and diffusion steps solved exactly; heat equation solved exactly}
\lv{Checks}{with 1024 modes and time step 0.005 the norms agree to a relative $6\cdot10^{-8}$; the energy balance $\|g(t)\|^2+2\nu\int_0^t\|(\partial_y g,g)\|^2ds=\|g(0)\|^2$ holds to $3\cdot10^{-7}$; decay rate of the norm on $[50,100]$: 0.0167, against $\nu^{1/2}\approx0.0316$}
\lv{Code}{\href{https://github.com/endipa/endipa.github.io/blob/main/numerics/heat-vs-shear/heat_vs_shear.py}{numerics/heat-vs-shear/heat\_vs\_shear.py} in the repository, also below (Python 3, NumPy, Matplotlib, ffmpeg); \texttt{python3 heat\_vs\_shear.py --out ../../wiki/media}, about a minute}
\lv{Status}{computed by Claude (AI) on 8 November 2026; not yet checked by a human \demo{}}
\end{levels}
\media{image}{heat-vs-shear-check.png}{Left: the norms of the two runs and of the heat equation. Right: the relative difference of the two runs and the energy-balance residual.}
\begin{note}[Code]
\begin{wikicode}[python]
#!/usr/bin/env python3
"""Heat equation versus advection-diffusion by a shear flow on the torus [0, 2pi)^2.

    d_t f + sin(y) d_x f = nu (d_xx + d_yy) f,      f(0, x, y) = cos x,

compared with the heat equation (the same without the transport term), whose solution is
exp(-nu t) cos x. The initial datum has zero average in x. It is a single Fourier mode in x,
so f(t, x, y) = Re(g(t, y) e^{ix}) with

    d_t g = -i sin(y) g + nu (d_yy - 1) g,      g(0, y) = 1,

solved by a Fourier method in y with Strang splitting: the transport step is the exact
multiplication by exp(-i sin(y) dt), the diffusion step is exact in Fourier variables.
The L2 norm reported is ||f(t)|| / ||f(0)||, that is the root mean square of |g(t, .)|.

Checks (written to heat-vs-shear-check.txt and heat-vs-shear-check.png):
  * the run is repeated with twice the Fourier modes and half the time step, and the norms are compared;
  * energy balance: ||g(t)||^2 + 2 nu int_0^t ||(d_y g, g)||^2 ds = ||g(0)||^2 (means over y).

The video lasts --duration seconds; frame s in [0, 1] shows t = T (a s + (1 - a) s^2), a = --slow (default 1:
time runs uniformly; a < 1 slows down the start); the last frame is held --hold seconds.

Outputs, in --out: heat-vs-shear.mp4 (video), heat-vs-shear-poster.png (one frame of it),
heat-vs-shear-check.png and heat-vs-shear-check.txt. Needs NumPy, Matplotlib and ffmpeg
(with the libx264 encoder). Run:  python3 heat_vs_shear.py --out ../../wiki/media

Drafted by Claude (AI) on 8 November 2026 for the wiki "Enhanced dissipation for passive
scalars"; record: page Numerical methods, section num-heat-vs-shear.
"""
import argparse
import os
import shutil
import subprocess
import tempfile

import numpy as np
import matplotlib

matplotlib.use("Agg")
import matplotlib.pyplot as plt

HEAT_COLOR = "#b4532a"
SHEAR_COLOR = "#1f4e79"


def solve(nu, T, N, dt, frame_times=()):
    """Times, norms ||f(t)||/||f(0)||, energy-balance residuals, and g(t, y) at frame_times."""
    y = 2 * np.pi * np.arange(N) / N
    eta = np.fft.fftfreq(N, d=1.0 / N)              # integer wavenumbers in y
    symbol = eta ** 2 + 1.0                          # symbol of -(d_yy - 1): mode k = 1 in x
    half = np.exp(-nu * symbol * dt / 2)
    transport = np.exp(-1j * np.sin(y) * dt)
    ghat = np.fft.fft(np.ones(N, dtype=complex))
    steps = int(round(T / dt))
    times = dt * np.arange(steps + 1)
    sq = np.empty(steps + 1)                         # mean over y of |g|^2
    dissip = np.empty(steps + 1)                     # mean over y of |d_y g|^2 + |g|^2
    frames, targets, j = [], list(frame_times), 0
    for n in range(steps + 1):
        power = np.abs(ghat) ** 2 / N ** 2           # Parseval: means over y
        sq[n] = power.sum()
        dissip[n] = (symbol * power).sum()
        while j < len(targets) and targets[j] <= times[n] + dt / 2:
            frames.append(np.fft.ifft(ghat))
            j += 1
        if n == steps:
            break
        ghat *= half
        ghat = np.fft.fft(np.fft.ifft(ghat) * transport) * half
    integral = np.concatenate([[0.0], np.cumsum((dissip[1:] + dissip[:-1]) / 2) * dt])
    residual = np.abs(sq + 2 * nu * integral - 1.0)
    return times, np.sqrt(sq), residual, frames, y


def field(g, nx):
    """f(x, y) = Re(g(y) e^{ix}) on an nx-by-nx grid (rows: y, columns: x)."""
    step = max(1, len(g) // nx)
    x = 2 * np.pi * np.arange(nx) / nx
    return np.real(g[::step][:, None] * np.exp(1j * x)[None, :])


def draw_frame(path, t, g, nu, times, norms, heat, nx, ymin):
    x = 2 * np.pi * np.arange(nx) / nx
    f_heat = np.exp(-nu * t) * np.cos(x)[None, :] * np.ones((nx, 1))
    f_shear = field(g, nx)
    fig = plt.figure(figsize=(11.2, 7.2), dpi=100)
    gs = fig.add_gridspec(2, 2, height_ratios=[3.6, 1.15], hspace=0.3, wspace=0.06,
                          left=0.08, right=0.97, top=0.845, bottom=0.095)
    for col, (data, title) in enumerate([(f_heat, "Heat equation:  $u = 0$"),
                                         (f_shear, r"Shear flow:  $u = (\sin y,\, 0)$")]):
        ax = fig.add_subplot(gs[0, col])
        ax.imshow(data, origin="lower", extent=[0, 2 * np.pi, 0, 2 * np.pi], cmap="RdBu_r",
                  vmin=-1, vmax=1, interpolation="bilinear")
        ax.set_title(title, fontsize=17, pad=8)
        ax.set_xticks([]); ax.set_yticks([])
        ax.set_aspect("equal")
    fig.suptitle(r"Same initial datum $f_0 = \cos x$, same diffusivity $\nu = 10^{-3}$"
                 + "        t = %d" % round(t), fontsize=17, y=0.975)
    ax = fig.add_subplot(gs[1, :])
    k = np.searchsorted(times, t, side="right")
    ax.semilogy(times, heat, color=HEAT_COLOR, alpha=0.18, lw=2)
    ax.semilogy(times, norms, color=SHEAR_COLOR, alpha=0.18, lw=2)
    ax.semilogy(times[:k], heat[:k], color=HEAT_COLOR, lw=2.4, label="heat equation")
    ax.semilogy(times[:k], norms[:k], color=SHEAR_COLOR, lw=2.4, label="shear flow")
    ax.axvline(t, color="#888888", lw=1)
    ax.set_xlim(0, times[-1]); ax.set_ylim(ymin, 1.5)
    ax.set_yticks([10.0 ** e for e in range(int(np.log10(ymin)), 1)])
    ax.set_xlabel("time  t", fontsize=13)
    ax.set_ylabel(r"$\|f(t)\|\,/\,\|f_0\|$", fontsize=14)
    ax.tick_params(labelsize=11)
    ax.legend(loc="lower left", fontsize=12, frameon=False, ncol=2)
    fig.savefig(path)
    plt.close(fig)


def main():
    ap = argparse.ArgumentParser(description=__doc__.split("\n")[0])
    ap.add_argument("--out", default=".", help="folder for the video and the figures")
    ap.add_argument("--nu", type=float, default=1e-3)
    ap.add_argument("--T", type=float, default=100.0)
    ap.add_argument("--N", type=int, default=512, help="Fourier modes in y")
    ap.add_argument("--dt", type=float, default=0.01)
    ap.add_argument("--duration", type=float, default=20.0, help="length of the video in seconds")
    ap.add_argument("--fps", type=int, default=24)
    ap.add_argument("--slow", type=float, default=1.0, help="speed of time at the start, relative to the mean (1: uniform)")
    ap.add_argument("--hold", type=float, default=1.5, help="seconds the last frame is held")
    ap.add_argument("--nx", type=int, default=256, help="pixels per side of each panel")
    a = ap.parse_args()
    os.makedirs(a.out, exist_ok=True)
    s = np.linspace(0, 1, int(round(a.duration * a.fps)))
    frame_times = a.T * (a.slow * s + (1 - a.slow) * s ** 2)
    times, norms, residual, frames, y = solve(a.nu, a.T, a.N, a.dt, frame_times)
    heat = np.exp(-a.nu * times)
    # check: the same run with N doubled and dt halved
    times2, norms2, residual2, _, _ = solve(a.nu, a.T, 2 * a.N, a.dt / 2)
    diff = np.abs(norms2[::2] - norms) / norms2[::2]
    ymin = 10 ** np.floor(np.log10(norms.min()))
    with open(os.path.join(a.out, "heat-vs-shear-check.txt"), "w") as fh:
        fh.write("nu=%g T=%g N=%d dt=%g frames=%d fps=%d\n" % (a.nu, a.T, a.N, a.dt, len(frame_times), a.fps))
        fh.write("final norm: shear %.4e, heat %.4e\n" % (norms[-1], heat[-1]))
        fh.write("max relative difference with (2N, dt/2): %.2e\n" % diff.max())
        fh.write("max energy-balance residual: %.2e (2N, dt/2: %.2e)\n" % (residual.max(), residual2.max()))
        fh.write("decay rate of the shear norm on [T/2, T]: %.4f; nu^(1/2) = %.4f\n" % (
            np.log(norms[len(norms) // 2] / norms[-1]) / (times[-1] - times[len(times) // 2]), np.sqrt(a.nu)))
    # check figure
    fig, (ax1, ax2) = plt.subplots(1, 2, figsize=(11, 4), dpi=100)
    ax1.semilogy(times, heat, color=HEAT_COLOR, lw=2, label="heat equation (exact)")
    ax1.semilogy(times, norms, color=SHEAR_COLOR, lw=2, label="shear flow, N=%d, dt=%g" % (a.N, a.dt))
    ax1.semilogy(times2, norms2, color="k", lw=1, ls="--", label="shear flow, N=%d, dt=%g" % (2 * a.N, a.dt / 2))
    ax1.set_xlabel("time t"); ax1.set_ylabel("norm ratio"); ax1.legend(frameon=False, fontsize=9)
    ax1.set_title("Norms (log scale)")
    ax2.semilogy(times[1:], diff[1:] + 1e-16, color=SHEAR_COLOR, lw=1.5, label="relative difference of the two runs")
    ax2.semilogy(times[1:], residual[1:] + 1e-16, color="#6b7280", lw=1.5, label="energy-balance residual")
    ax2.set_xlabel("time t"); ax2.legend(frameon=False, fontsize=9); ax2.set_title("Checks")
    fig.tight_layout()
    fig.savefig(os.path.join(a.out, "heat-vs-shear-check.png"))
    plt.close(fig)
    # video
    with tempfile.TemporaryDirectory() as tmp:
        for i, (t, g) in enumerate(zip(frame_times, frames)):
            draw_frame(os.path.join(tmp, "f%04d.png" % i), t, g, a.nu, times, norms, heat, a.nx, ymin)
        poster = min(range(len(frame_times)), key=lambda i: abs(frame_times[i] - 0.2 * a.T))
        shutil.copyfile(os.path.join(tmp, "f%04d.png" % poster), os.path.join(a.out, "heat-vs-shear-poster.png"))
        subprocess.run(["ffmpeg", "-y", "-loglevel", "error", "-framerate", str(a.fps),
                        "-i", os.path.join(tmp, "f%04d.png"), "-vf", "tpad=stop_mode=clone:stop_duration=%g" % a.hold,
                        "-c:v", "libx264", "-pix_fmt", "yuv420p",
                        "-crf", "26", "-preset", "slow", "-movflags", "+faststart",
                        os.path.join(a.out, "heat-vs-shear.mp4")], check=True)
    print(open(os.path.join(a.out, "heat-vs-shear-check.txt")).read().strip())


if __name__ == "__main__":
    main()
\end{wikicode}
\end{note}

\subsection{Alternating shears with random phases}\label{num-alternating-shears}
\begin{levels}
\lv{Produces}{the second video on the \wikilink{Overview}{} (\texttt{alternating-shears.mp4}) and the check figure below}
\lv{Equation}{$\partial_t f+u\cdot\nabla f=\nu\Delta f$ on $[0,2\pi)^2$, periodic, with $f_0=\cos x$ and the alternating shears of Pierrehumbert's model: $u=(\sin(y+\varphi_n),0)$ for $6n\le t<6n+3$ and $u=(0,\sin(x+\psi_n))$ for $6n+3\le t<6n+6$, the phases independent and uniform on $[0,2\pi)$; compared with the steady shear $u=(\sin y,0)$ and with the heat equation. Both velocity fields have size at most 1.}
\lv{Parameters}{$\nu=10^{-4}$, each shear acting for a time 3, times $0\le t\le 60$, phases from the random generator of NumPy with seed 2026; video of 20 s at 24 frames per second, time running uniformly (3 time units per second), last frame held 1.5 s}
\lv{Method}{Fourier pseudo-spectral method on a $512\times512$ grid; Strang splitting with time step 0.0125, the transport steps (each row or column shifted, through its Fourier coefficients) and the diffusion steps solved exactly; the steady shear as in the record above}
\lv{Checks}{with $1024\times1024$ modes and time step 0.00625 the norms agree to a relative $1.2\cdot10^{-6}$; the energy balance $\|f(t)\|^2+2\nu\int_0^t\|\nabla f\|^2ds=\|f_0\|^2$ holds to $3.5\cdot10^{-6}$; time for the norm to halve, for $\nu=10^{-2},10^{-3},10^{-4},10^{-5}$: alternating shears 8.8, 16.3, 23.9, 31.7 (the same with twice the modes), an increase of about 7.6 for each factor 10 in $\nu$; steady shear 8.6, 20.0, 43.9, 94.9; heat equation $\ln 2/\nu$}
\lv{Context}{Exponential mixing for randomly shifted alternating shears is proved by \citet{BCZG2023}, and dissipation on the time scale $|\log\nu|$, with high probability, by \citet{CoopermanIyerSon2025}; whether the parameters used here satisfy their hypotheses has not been checked}
\lv{Code}{\href{https://github.com/endipa/endipa.github.io/blob/main/numerics/alternating-shears/alternating_shears.py}{numerics/alternating-shears/alternating\_shears.py} in the repository (Python 3, NumPy, SciPy, Matplotlib, ffmpeg); \texttt{python3 alternating\_shears.py --out ../../wiki/media}, about 15 minutes}
\lv{Status}{computed by Claude (AI) on 8 November 2026; not yet checked by a human \demo{}}
\end{levels}
\media{image}{alternating-shears-check.png}{Left: the norms for $\nu=10^{-4}$, with two resolutions for the alternating shears. Right: the time for the norm to halve, against $1/\nu$, logarithmic scales.}

\section{Solvers}
\placeholder{The methods used across the records, with their strengths and limits.}

\section*{References}
%%wiki <!-- refs-begin: BCZG2023 CoopermanIyerSon2025 -->
\begin{thebibliography}{}
\bibitem[Blumenthal, Coti Zelati and Gvalani(2023)]{BCZG2023} A. Blumenthal, M. Coti Zelati and R. S. Gvalani, \emph{Exponential mixing for random dynamical systems and an example of Pierrehumbert}, Annals of Probability \textbf{51} (2023), 1559–1601, \href{https://arxiv.org/abs/2204.13651}{arXiv:2204.13651}.
\bibitem[Cooperman, Iyer and Son(2025)]{CoopermanIyerSon2025} W. Cooperman, G. Iyer and S. Son, \emph{A Harris theorem for enhanced dissipation, and an example of Pierrehumbert}, Nonlinearity \textbf{38} (2025), 045027, \href{https://arxiv.org/abs/2403.19858}{arXiv:2403.19858}.
\end{thebibliography}
%%wiki <!-- refs-end -->
\end{document}
